\chapter{Inleiding}
\label{sec:inleiding}
Deze \MakeLowercase\documenttype{} laat zien hoe je de \LaTeX{}-stijl die is ontwikkeld voor documenten van de opleiding Elektrotechniek van de Hogeschool Rotterdam kunt gebruiken. Deze stijl is gebaseerd op een stijl die eerder door de auteur ontwikkeld is bij de opleiding Elektrotechniek aan De Haagse Hogeschool en op een Word-template van de Hogeschool Rotterdam die ontwikkeld is door Johan Peltenburg. 

Met deze template kan vanuit één \LaTeX{}-bestand zowel een papieren als een ebook versie gegenereerd worden. De papieren versie is bedoeld om dubbelzijdig te printen, heeft een bindmarge, hoofdstukken beginnen op oneven pagina's en links hebben geen kleur (klikken op papier werkt namelijk niet, de links werken overigens wel als het document toch op een elektronisch device wordt gelezen). De ebook versie is enkelzijdig en heeft gekleurde links om de lezer er op te attenderen dat je er op kunt klikken. In de rest van deze \MakeLowercase\documenttype{} wordt uitgelegd hoe je deze \LaTeX{}-stijl kunt gebruiken en worden voorbeelden gegeven om te laten zien hoe de layout van bepaalde dingen is.

Er is ook een versie van deze template beschikbaar voor het maken van Engelstalige documenten. Specifieke aandachtspunten voor het gebruik van die template worden in een apart hoofdstuk besproken. Tot slot volgt nog een hoofdstuk met enkele algemene tips voor het gebruik van \LaTeX{}. 

Deze stijl is getest met de \TeX{}-distributie Tex Live\footnote{\url{ http://www.tug.org/texlive/}} en de verwachting is dat deze stijl met andere distributies net zo goed werkt, maar dat is niet getest.

Zelf maak ik gebruik van de \LaTeX{}-editor TeXstudio\footnote{\url{http://texstudio.sourceforge.net/}}. Er is een zogenoemd session-
bestand beschikbaar waarmee in één keer alle relevante bestanden geopend worden: \code{main.txss}. Maar je kunt deze template ook prima met elke andere editor gebruiken.

Voorbeelden van documenten (book style) die met deze template gemaakt zijn:
\begin{itemize}
\item \citetitle{dcs01} \cite{dcs01} (Nederlands).\footnote{Sourcecode: \url{http://www.medewerker.hro.nl/brojz/DCS01/Collegedictaat_DCS01.zip}.}
\item \citetitle{tds02} \cite{tds02} (Engels).\footnote{Sourcecode: \url{http://www.medewerker.hro.nl/brojz/TDS02/Lab_Work_Handbook_TDS02.zip}.}
\item Dit document zelf.\footnote{Sourcecode: \url{http://www.medewerker.hro.nl/brojz/Handleiding_LaTeX-stijl_HR.zip}.}
\end{itemize} 

Voorbeelden van documenten (article style) die met deze template gemaakt zijn:
\begin{itemize}
\item Modulewijzer TDS02 (Nederlands).\footnote{Zie: \url{http://www.medewerker.hro.nl/brojz/TDS02/Modulewijzer_TDS02.pdf} en \url{http://www.medewerker.hro.nl/brojz/TDS02/Modulewijzer_TDS02_ebook.pdf}. De sourcecode vind je hier: \url{http://www.medewerker.hro.nl/brojz/TDS02/Modulewijzer_TDS02.zip}.}
\item Report Requirements for TDS02 (Engels).\footnote{Zie: \url{http://www.medewerker.hro.nl/brojz/TDS02/Report_Requirements_TDS02.pdf} en \url{http://www.medewerker.hro.nl/brojz/TDS02/Report_Requirements_TDS02_ebook.pdf}. De sourcecode vind je hier: \url{http://www.medewerker.hro.nl/brojz/TDS02/Report_Requirements_TDS02.zip}.}
\end{itemize} 

Op- en aanmerkingen zijn altijd welkom. Je kunt me mailen op \url{mailto:J.Z.M.Broeders@hr.nl}.
